\documentclass[11pt]{article}

\usepackage[T1]{fontenc}
\usepackage[utf8]{inputenc}
\usepackage{lmodern}
\usepackage{microtype}
\usepackage[a4paper,margin=1.05in]{geometry}
\usepackage{amsmath,amssymb,amsthm,mathtools}
\usepackage{booktabs,longtable,array}
\usepackage{enumitem}
\usepackage{xcolor}
\usepackage{hyperref}
\usepackage[nameinlink,capitalise,noabbrev]{cleveref}
\usepackage{url}
\usepackage{listings}
\usepackage{fancyhdr}
\usepackage{tikz}
\usetikzlibrary{arrows.meta,positioning}

\definecolor{deepblue}{RGB}{29,64,112}
\definecolor{deepred}{RGB}{137,42,42}
\definecolor{softgray}{RGB}{245,246,248}
\hypersetup{
  colorlinks=true,
  linkcolor=deepblue,
  citecolor=deepblue,
  urlcolor=deepblue,
  pdftitle={Nondeterministic Width-One Algebraic Branching Programs in Characteristic Two},
  pdfauthor={Author information to be supplied}
}

\pagestyle{fancy}
\fancyhf{}
\fancyhead[L]{Nondeterministic width-one ABPs in characteristic two}
\fancyhead[R]{Preprint draft}
\fancyfoot[C]{\thepage}
\setlength{\headheight}{14pt}

\newtheorem{theorem}{Theorem}[section]
\newtheorem{lemma}[theorem]{Lemma}
\newtheorem{proposition}[theorem]{Proposition}
\newtheorem{corollary}[theorem]{Corollary}
\newtheorem{claim}[theorem]{Claim}
\theoremstyle{definition}
\newtheorem{definition}[theorem]{Definition}
\newtheorem{construction}[theorem]{Construction}
\newtheorem{remark}[theorem]{Remark}
\newtheorem{example}[theorem]{Example}

\newcommand{\F}{\mathbb{F}}
\newcommand{\B}{\{0,1\}}
\newcommand{\Ncl}{\mathsf{N}}
\newcommand{\VP}{\mathbf{VP}}
\newcommand{\VNP}{\mathbf{VNP}}
\newcommand{\VPe}{\mathbf{VP}_{\!\mathrm e}}
\newcommand{\VNPe}{\mathbf{VNP}_{\!\mathrm e}}
\newcommand{\VPone}{\mathbf{VP}_{1}}
\newcommand{\VNPone}{\mathbf{VNP}_{1}}
\newcommand{\VNPtwo}{\mathbf{VNP}_{1}^{[\le 2]}}
\newcommand{\VNPwplus}{\mathbf{VNP}_{1}^{\mathrm{w}+}}
\newcommand{\supp}{\operatorname{supp}}
\newcommand{\val}{\operatorname{val}}
\newcommand{\indeg}{\operatorname{indeg}}
\newcommand{\outdeg}{\operatorname{outdeg}}
\newcommand{\poly}{\operatorname{poly}}
\newcommand{\charac}{\operatorname{char}}
\newcommand{\one}{\mathbf{1}}
\newcommand{\lean}{\textsc{Lean}}
\newcommand{\code}[1]{\texttt{\detokenize{#1}}}

\DeclareMathOperator{\Eval}{Eval}
\DeclareMathOperator{\Flow}{Flow}
\DeclareMathOperator{\Sel}{Sel}
\DeclareMathOperator{\Pairs}{Pairs}
\DeclareMathOperator{\td}{tdeg}

\lstdefinestyle{leanstyle}{
  basicstyle=\ttfamily\small,
  columns=fullflexible,
  breaklines=true,
  frame=single,
  backgroundcolor=\color{softgray},
  rulecolor=\color{black!15},
  xleftmargin=0.5em,
  xrightmargin=0.5em
}
\lstset{style=leanstyle}

\title{\bfseries Nondeterministic Width-One Algebraic Branching Programs\\
       in Characteristic Two}
\author{Author information to be supplied\\
\small Preprint draft accompanying a machine-checked \lean\ artifact}
\date{19 September 2026}

\begin{document}
\maketitle

\begin{center}
\fcolorbox{deepred}{softgray}{\parbox{0.91\textwidth}{
\textbf{Draft status.}
The mathematical core described below is represented in a \lean\ 4 archive and the
accompanying artifact report records a successful build and axiom audit.  The public
literature consequences use the standard identifications of Valiant and
Bringmann--Ikenmeyer--Zuiddam and are not imported as axioms.  This manuscript has not yet
undergone external peer review; a final bibliographic novelty audit should precede public
claims of priority.
}}
\end{center}

\begin{abstract}
Bringmann, Ikenmeyer and Zuiddam proved that nondeterministic width-one algebraic
branching programs capture \(\VNP\) over fields of characteristic different from two,
and, more strongly, that every edge label may involve at most two variables.  They also
proved a strict separation over \(\F_2\), leaving the other fields of characteristic two
open.  We give a characteristic-two compiler that covers every field
\(\F\neq\F_2\).

Let \(\Phi\) be a binary division-free arithmetic formula of size \(s\) over a field
\(\F\) of characteristic two.  If \(\tau\in\F\setminus\{0,1\}\), then \(\Phi\) admits a
representation
\[
  \Phi(\mathbf{x})
   = \sum_{\mathbf{b}\in\B^q}\prod_{j=1}^{M} L_j(\mathbf{x},\mathbf{b}),
\]
where every \(L_j\) is affine and involves at most two variables, with
\(q\le 14s\) and \(M\le 34s\).  The construction rests on a one-bit, three-factor
identity for \(1+xy\), a one-bit ternary identity with an explicit cubic error, and a
path-selector polynomial whose existing pair-exclusion factors kill that error only after
the selector variables have been fixed to Boolean values.  The proof order is essential:
the corresponding cancellation is not a formal polynomial identity.

At the class level we obtain
\[
  \Ncl(\VPe(\F))=\VNPtwo(\F)
  \qquad
  (\charac(\F)=2,\ \F\neq\F_2).
\]
Combined with Valiant's \(\VNPe=\VNP\), the results of
Bringmann--Ikenmeyer--Zuiddam, and their separation over \(\F_2\), this yields the complete
field classification
\[
  \VNPwplus(\F)=\VNP(\F) \iff |\F|>2,
  \qquad
  \VNPone(\F)=\VNP(\F) \iff \F\not\cong\F_2.
\]
The finite compiler, its class lift, its width-one literature bridge, and the relevant
scope firewalls have been formalized in \lean\ 4.
\end{abstract}

\tableofcontents


\section{Introduction}

Algebraic branching programs (ABPs) interpolate between formulas and general arithmetic
circuits.  A width-one ABP is simply a directed path, and hence computes a product of
affine linear forms.  Applying a Boolean hypercube summation to such products gives the
nondeterministic class \(\VNPone\).  Bringmann, Ikenmeyer and Zuiddam
\cite{BringmannIkenmeyerZuiddam2018,BringmannIkenmeyerZuiddam2017} proved
\[
  \VNPone(\F)=\VNP(\F)
  \quad\text{when }\charac(\F)\neq 2,
\]
and sharpened their construction so that each affine factor uses at most two variables.
Their argument uses identities containing division by two.  They showed separately that
\(\VNPone(\F_2)\subsetneq\VNP(\F_2)\), and explicitly left the remaining
characteristic-two fields open.

This paper supplies the missing characteristic-two construction whenever the field has an
element other than zero and one.  The exceptional field is therefore exactly \(\F_2\), in
agreement with the earlier separation.

\subsection{Main finite theorem}

We count the size of a binary arithmetic formula by the number of leaves, addition gates,
and multiplication gates.  Leaves are variables or field constants.

\begin{theorem}[Support-two characteristic-two compiler]\label{thm:finite-main}
Let \(\F\) be a field of characteristic two, let
\(\tau\in\F\setminus\{0,1\}\), and let \(\Phi\) be a binary division-free arithmetic
formula of size \(s\).  There exist integers \(q,M\ge 0\) and affine linear forms
\[
  L_1,\ldots,L_M\in
  \F[\mathbf{x},b_1,\ldots,b_q]
\]
such that each \(L_j\) involves at most two variables and
\[
  \Phi(\mathbf{x})
  =\sum_{\mathbf{b}\in\B^q}\prod_{j=1}^{M}L_j(\mathbf{x},\mathbf{b}).
\]
Moreover,
\[
  q\le 14s,
  \qquad
  M\le 34s.
\]
\end{theorem}

The formal artifact proves a stronger syntactic statement for its generated factors: each
factor contains at most two variable terms in its explicit affine-form term list.  This
implies the usual semantic support bound appearing in \cref{thm:finite-main}.

\subsection{Class consequences}

Write \(\VNPtwo(\F)\) for polynomial families admitting polynomial-dimensional Boolean
hypercube sums of polynomially many affine forms, each involving at most two variables.
Write \(\Ncl(\mathcal C)\) for the nondeterministic closure of a class \(\mathcal C\).

\begin{theorem}[Internal class equality]\label{thm:class-main}
Let \(\F\) have characteristic two and suppose \(\F\neq\F_2\).  Then
\[
   \Ncl(\VPe(\F))=\VNPtwo(\F).
\]
The statement covers finite and infinite fields.
\end{theorem}

A product of affine forms is a width-one ABP: place the factors consecutively on a directed
path.  Thus the support-two class maps directly into the class denoted
\(\VNPwplus\) by Bringmann--Ikenmeyer--Zuiddam.  Conversely, constant-width ABPs have
polynomial-size formulas, and Valiant's characterization gives
\(\Ncl(\VPe)=\VNP\) \cite{Valiant1980,BringmannIkenmeyerZuiddam2018}.

\begin{corollary}[Complete field classification]\label{cor:classification}
For every field \(\F\),
\[
  \VNPwplus(\F)=\VNP(\F)
  \quad\Longleftrightarrow\quad |\F|>2.
\]
For unrestricted affine labels,
\[
  \VNPone(\F)=\VNP(\F)
  \quad\Longleftrightarrow\quad \F\not\cong\F_2.
\]
\end{corollary}

For characteristic different from two, the first equality is the support-two theorem of
Bringmann--Ikenmeyer--Zuiddam.  For characteristic two and \(\F\neq\F_2\), it follows from
\cref{thm:class-main}.  Over \(\F_2\), even unrestricted width one is strictly weaker than
\(\VNP\) by their Proposition~9.1.

\subsection{Proof architecture}

The proof is deliberately modular.

\begin{enumerate}[label=\textbf{Step \arabic*.},leftmargin=2.2cm]
\item Convert a formula into a ranked, labelled acyclic graph whose path polynomial is the
formula.  The graph has maximum total degree three and labels involving at most one
original variable.
\item Introduce one Boolean selector variable for each edge.  A product of affine flow
factors and quadratic pair exclusions is the indicator of a selected source--sink path.
Multiplying by conditional edge-label factors and summing over all selectors recovers the
ABP path polynomial.
\item Replace every quadratic factor \(1+xy\) by a one-bit sum of three support-two affine
factors.
\item Replace each ternary internal flow factor by a one-bit sum of four support-two affine
factors.  The new identity carries an explicit cubic error.
\item Fix the original edge-selector assignment.  At every ternary vertex, two incident
edges share a head or a tail, so an already-present pair exclusion kills the cubic error
pointwise on the Boolean cube.  Only after this cancellation do we sum over the original
selectors.
\item Count the fresh auxiliaries and factors.  A sharp structural identity shows that the
number of ordered pair exclusions is four times the number of addition gates.
\item Flatten an outer \(\VNP\)-type witness cube with the compiler's fresh cube to obtain
the class-level inclusion.  The reverse inclusion follows because a product of
polynomially many affine forms has a polynomial-size formula.
\end{enumerate}

The distinction in Step~5 is not cosmetic.  The polynomial
\(abc(1+ab)\) is nonzero in \(\F[a,b,c]\); it vanishes only after \(a,b,c\) are specialized
to Boolean values.  The formal development contains an explicit negative theorem
preventing this pointwise argument from being silently promoted to a polynomial identity.

\subsection{Scope}

The result is a representation theorem.  The hypercube may have linearly many witness
bits, and no faster evaluator for the resulting sum is provided.  Consequently, nothing
here separates \(\VP\) from \(\VNP\), resolves a Boolean complexity-class problem, or
implies a subexponential algorithm.

The characteristic-two border-width-two problem studied by Dutta et al.
\cite{DuttaEtAl2024} is adjacent but different.  Their model is deterministic and
border-approximate; the present construction is exact and nondeterministic.  No
Frobenius-root extraction or border simulation is claimed here.

\section{Preliminaries}\label{sec:prelim}

\subsection{Polynomials, formulas, and families}

Throughout, \(\F\) is a field.  A \emph{binary arithmetic formula} is a rooted binary tree
whose internal nodes are labelled by \(+\) or \(\times\), and whose leaves are variables
or elements of \(\F\).  Its size is the number of nodes.  If \(l,a,u\) denote the numbers
of leaves, addition gates, and multiplication gates, respectively, then
\[
  s=l+a+u.
\]
A family \(f=(f_n)_{n\ge 0}\) is a \emph{p-family} if its number of variables and total
degree are polynomially bounded.  We use standard definitions of \(\VPe\) and \(\VNP\);
see, for example, \cite{Burgisser2000,MalodPortier2008}.  The class \(\VPe\) consists of
p-families with polynomially bounded formula size.

\subsection{Algebraic branching programs}

An algebraic branching program is a finite directed acyclic graph \(G=(V,E)\) with a
source \(s\), a sink \(t\), and an affine label \(\ell_e\) on every edge.  Its value is
\[
  \val(G)=\sum_{P:s\leadsto t}\prod_{e\in P}\ell_e.
\]
A width-one ABP is a directed path and hence computes a product of affine forms.  We use
\(\VPone\) for polynomial-size width-one ABPs with unrestricted affine labels, and
\(\VPone^{\mathrm{w}+}\) for the version in which every edge label involves at most two
variables.  Following Bringmann--Ikenmeyer--Zuiddam, set
\[
  \VNPone=\Ncl(\VPone),
  \qquad
  \VNPwplus=\Ncl(\VPone^{\mathrm{w}+}).
\]

\subsection{Nondeterministic closure}

\begin{definition}[Nondeterministic closure]\label{def:nclosure}
For a class \(\mathcal C\) of polynomial families, \(\Ncl(\mathcal C)\) consists of all
families \(f=(f_n)\) for which there exist a family \(g=(g_n)\in\mathcal C\) and
polynomially bounded functions \(p,q\) such that
\[
  f_n(\mathbf{x})
   =\sum_{\mathbf{b}\in\B^{p(n)}}g_{q(n)}(\mathbf{b},\mathbf{x}).
\]
\end{definition}

This is the closure operator introduced as Definition~2.2 in
\cite{BringmannIkenmeyerZuiddam2018}.  We will use its monotonicity:
\[
  \mathcal C\subseteq\mathcal D
  \quad\Longrightarrow\quad
  \Ncl(\mathcal C)\subseteq\Ncl(\mathcal D).
\]

\subsection{Affine-product hypercube representations}

\begin{definition}[Support-two affine-hypercube representation]\label{def:rep}
Let \(f\in\F[\mathbf{x}]\).  A support-two affine-hypercube representation of \(f\) is an
identity
\[
  f(\mathbf{x})
  =\sum_{\mathbf{b}\in\B^q}\prod_{j=1}^{M}L_j(\mathbf{x},\mathbf{b}),
\]
where every \(L_j\) is affine and \(|\supp(L_j)|\le 2\).  Constants are allowed, and the
support counts both original and hypercube variables.
\end{definition}

For families, \(q\) and \(M\) must be polynomially bounded.  The resulting class is denoted
\(\VNPtwo\).  Since a product of \(M\) affine factors is computed by a width-one path of
length \(M\), one has the immediate bridge
\[
  \VNPtwo\subseteq\VNPwplus.
\]

\section{Characteristic-two gadgets}\label{sec:gadgets}

Assume throughout this section that \(\charac(\F)=2\).  Fix
\(\tau\in\F\setminus\{0,1\}\) and put
\[
  \kappa=\tau(\tau+1).
\]

\begin{lemma}[Denominator legality]\label{lem:kappa}
The element \(\kappa\) is nonzero.
\end{lemma}

\begin{proof}
Both \(\tau\) and \(\tau+1\) are nonzero.  Indeed, \(\tau+1=0\) would imply
\(\tau=-1=1\) in characteristic two.  Since \(\F\) is a field, their product is nonzero.
\end{proof}

\subsection{A one-bit quadratic gadget}

\begin{lemma}[Quadratic identity]\label{lem:quadratic}
For arbitrary elements \(x,y\) of any commutative \(\F\)-algebra,
\[
\sum_{h\in\B}
  \kappa^{-1}(\tau+1+h)(x+\tau+h)(\kappa y+\tau+h)
 =1+xy.
\]
Only \(h\) is Boolean; no Boolean relation is imposed on \(x\) or \(y\).
\end{lemma}

\begin{proof}
First prove the denominator-free identity.  The left-hand side before multiplication by
\(\kappa^{-1}\) is
\[
 S=(\tau+1)(x+\tau)(\kappa y+\tau)
   +\tau(x+\tau+1)(\kappa y+\tau+1).
\]
The coefficient of \(\kappa y\) is
\[
 (\tau+1)(x+\tau)+\tau(x+\tau+1)
 =x\bigl((\tau+1)+\tau\bigr)
   +2\tau(\tau+1)=x.
\]
The remaining term is
\[
 \tau(\tau+1)\bigl((x+\tau)+(x+\tau+1)\bigr)=\kappa.
\]
Thus \(S=\kappa xy+\kappa=\kappa(1+xy)\).  Divide by \(\kappa\), using
\cref{lem:kappa}.
\end{proof}

The three factors in each summand involve, respectively, only \(h\), the pair \((x,h)\),
and the pair \((y,h)\), whenever \(x\) and \(y\) are affine expressions supported on one
variable each.  A distinct fresh \(h\) will be used for every occurrence of the gadget.

\subsection{A ternary gadget with a controlled error}

Define the two-point functional
\[
 \Lambda_\tau[p]
   =(\tau+1)p(\tau)+\tau p(\tau+1).
\]
A direct calculation in characteristic two gives
\[
 \Lambda_\tau[1]=1,
 \qquad
 \Lambda_\tau[T]=0,
 \qquad
 \Lambda_\tau[T^2]=\kappa,
 \qquad
 \Lambda_\tau[T^3]=\kappa.
\]

\begin{lemma}[Ternary identity]\label{lem:ternary}
For arbitrary elements \(a,b,c\) of any commutative \(\F\)-algebra,
\begin{align*}
&\sum_{h\in\B}
 \kappa^{-1}(\tau+1+h)(a+\tau+h)(b+\tau+h)(c+\tau+h)\\
&\hspace{35mm}=1+a+b+c+\kappa^{-1}abc.
\end{align*}
\end{lemma}

\begin{proof}
Apply \(\Lambda_\tau\) to
\[
 p(T)=(a+T)(b+T)(c+T)
 =abc+(ab+ac+bc)T+(a+b+c)T^2+T^3.
\]
The four moments displayed above yield
\[
 \Lambda_\tau[p]=abc+\kappa(a+b+c)+\kappa
 =\kappa(1+a+b+c)+abc.
\]
Division by \(\kappa\) proves the claim.
\end{proof}

The right-hand side contains the genuine cubic error \(\kappa^{-1}abc\).  It must not be
discarded as a formal polynomial.

\begin{lemma}[Boolean error killer]\label{lem:killer}
For \(u,v,w\in\B\), viewed inside any characteristic-two field,
\[
  uvw(1+uv)=0.
\]
\end{lemma}

\begin{proof}
If \(uv=0\), the product vanishes.  If \(uv=1\), then \(1+uv=1+1=0\).
\end{proof}

\begin{remark}[Formal versus pointwise cancellation]\label{rem:firewall}
The polynomial
\[
 abc(1+ab)=abc+a^2b^2c
\]
is nonzero in \(\F[a,b,c]\).  The cancellation in \cref{lem:killer} is valid only after
specializing \(a,b,c\) to Boolean values.  This distinction determines the proof order in
\cref{sec:compiler} and is separately machine-checked in the formal artifact.
\end{remark}

\section{From formulas to bounded-degree labelled DAGs}\label{sec:formula-dag}

We recall a standard syntax-directed conversion.  Its precise form is useful because the
support and counting arguments depend on local degrees.

\begin{construction}[Formula DAG]\label{cons:dag}
Associate to each formula \(\Phi\) a labelled ranked DAG \(G_\Phi\) as follows.
\begin{enumerate}[label=(\roman*)]
\item A leaf \(\ell\) becomes one edge from a fresh source to a fresh sink, labelled
\(\ell\).
\item If \(\Phi=\Psi\Theta\), take the disjoint union of \(G_\Psi\) and \(G_\Theta\) and
add one unit-labelled edge from the sink of \(G_\Psi\) to the source of \(G_\Theta\).
\item If \(\Phi=\Psi+\Theta\), take the disjoint union, add a fresh source with one
unit-labelled edge to each child source, and add one unit-labelled edge from each child
sink to a fresh sink.
\end{enumerate}
Ranks are chosen compatibly so every edge increases rank.
\end{construction}

\begin{lemma}[Value and local degree]\label{lem:dag-properties}
The path polynomial of \(G_\Phi\) equals \(\Phi\).  Every edge label is a variable or a
constant.  The source has indegree zero and outdegree at most two; the sink has outdegree
zero and indegree at most two; every internal vertex has total degree at most three.
\end{lemma}

\begin{proof}
The value statement follows by structural induction.  A leaf has its label as its unique
path value.  The multiplication construction concatenates a path in the left graph with
the unit connector and a path in the right graph, multiplying their values.  The addition
construction offers a disjoint choice between the two child graphs, adding their path
sums.  The degree assertions follow simultaneously from the three constructors.  In the
multiplication constructor, the old left sink gains one outgoing connector and the old
right source gains one incoming connector; each had the opposite boundary degree zero.
In the addition constructor, the new boundary vertices have degree two, while each child
boundary gains one incident connector.  Thus no internal total degree exceeds three.
\end{proof}

Let \(l,a,u\) be the leaf, addition-gate, and multiplication-gate counts of \(\Phi\).

\begin{lemma}[Graph counts]\label{lem:graph-counts}
For \(G_\Phi=(V,E)\),
\[
  |E|=l+4a+u,
  \qquad
  |V|=2l+2a.
\]
\end{lemma}

\begin{proof}
Induct on the formula.  A leaf contributes two vertices and one edge.  A multiplication
adds the counts of the children and one connector edge.  An addition adds the child counts,
two new vertices, and four connector edges.
\end{proof}

\section{The Boolean path selector}\label{sec:selector}

Fix a ranked DAG \(G=(V,E)\) satisfying the degree conditions of
\cref{lem:dag-properties}.  Introduce a Boolean variable \(z_e\) for each edge.  For a
vertex \(v\), set
\[
 I_v(z)=\sum_{e:\,\mathrm{tgt}(e)=v}z_e,
 \qquad
 O_v(z)=\sum_{e:\,\mathrm{src}(e)=v}z_e.
\]
Define the flow factor
\[
 F_v(z)=
 \begin{cases}
   O_s(z),&v=s,\\
   I_t(z),&v=t,\\
   1+I_v(z)+O_v(z),&v\notin\{s,t\}.
 \end{cases}
\]
For every ordered pair of distinct edges with a common head, and every ordered pair of
distinct edges with a common tail, include a pair-exclusion factor \(1+z_ez_{e'}\).  Let
\(\mathcal P\) be the resulting ordered-pair index set, and put
\[
 \Sel_G(z)=
 \left(\prod_{v\in V}F_v(z)\right)
 \left(\prod_{(e,e')\in\mathcal P}(1+z_ez_{e'})\right).
\]
The use of ordered pairs duplicates each unordered exclusion.  This does not change its
Boolean value, which is zero or one.

\begin{proposition}[Selector correctness]\label{prop:selector}
For every \(z\in\B^E\),
\[
 \Sel_G(z)=
 \begin{cases}
  1,&\{e:z_e=1\}\text{ is the edge set of a directed }s\text{--}t\text{ path},\\
  0,&\text{otherwise}.
 \end{cases}
\]
\end{proposition}

\begin{proof}
Every factor evaluates to zero or one.  The pair exclusions force at most one selected
incoming edge and at most one selected outgoing edge at each vertex.  Subject to those
exclusions, the source flow factor is one precisely when the source has one selected
outgoing edge, and the sink factor is one precisely when the sink has one selected incoming
edge.  At an internal vertex, \(I_v,O_v\in\{0,1\}\), and
\(1+I_v+O_v=1\) precisely when \(I_v=O_v\).  Hence every selected internal vertex has one
selected incoming edge if and only if it has one selected outgoing edge.

Starting at the unique selected edge out of \(s\), conservation and finiteness extend it
forward until it reaches \(t\).  It cannot stop at an internal vertex.  It cannot revisit a
vertex because the graph is ranked and acyclic.  Thus it gives an \(s\)-\(t\) path.  Any
additional selected component would be finite, avoid \(s,t\), and have equal selected
indegree and outdegree at every one of its vertices.  Following its unique outgoing edges
would create a directed cycle, contradicting acyclicity.  Therefore no additional component
exists.  The converse is immediate: an \(s\)-\(t\) path satisfies all flow and exclusion
factors.
\end{proof}

Let \(\ell_e\) be the affine label of edge \(e\), and define its conditional label factor
\[
  C_e(z_e)=1+z_e(\ell_e+1).
\]
In characteristic two,
\[
  C_e(0)=1,
  \qquad
  C_e(1)=\ell_e.
\]

\begin{proposition}[ABP hypercube identity]\label{prop:abp-cube}
For a labelled DAG over a field of characteristic two,
\[
  \val(G)=
  \sum_{z\in\B^E}\Sel_G(z)\prod_{e\in E}C_e(z_e).
\]
\end{proposition}

\begin{proof}
By \cref{prop:selector}, only selector assignments that are edge sets of source--sink paths
survive.  For such an assignment, the conditional label product equals the product of the
labels on the selected path; all unselected edges contribute one.  Summing the surviving
terms is exactly the path polynomial.
\end{proof}

\section{The support-two compiler}\label{sec:compiler}

The expression in \cref{prop:abp-cube} is not yet a product of support-two affine forms.
It contains quadratic factors \(1+xy\) and, at total-degree-three internal vertices,
ternary affine flow factors.  We now compile both kinds with fresh one-bit gadgets.

\subsection{Occurrence sets and fresh auxiliaries}

Let
\[
  D=\#\{v\in V\setminus\{s,t\}:\indeg(v)+\outdeg(v)=3\}
\]
and let \(P=|\mathcal P|\) be the ordered pair count.  There are exactly
\(|E|+P\) quadratic occurrences: one conditional edge factor per edge and one exclusion
per ordered pair.  Give every such occurrence its own fresh bit.  Give every degree-three
internal flow factor its own fresh bit.  The full auxiliary set is the disjoint union of
\[
  \underbrace{E}_{\text{original selectors}}
  \ \sqcup\ 
  \underbrace{T}_{\text{ternary gadgets}}
  \ \sqcup\ 
  \underbrace{K}_{\text{quadratic gadgets}},
\]
where \(|T|=D\) and \(|K|=|E|+P\).

Freshness is structural.  It prevents the generally false replacement
\[
  \left(\sum_h A_h\right)\left(\sum_k B_k\right)
  \stackrel{\text{false}}{=}
  \sum_h A_hB_h.
\]
With disjoint bits, finite distributivity gives a sum over the Cartesian product of the
local cubes.

\subsection{Compiling quadratic occurrences}

Every pair exclusion is \(1+z_ez_{e'}\).  Every conditional edge factor can be written
\[
  1+z_e(\ell_e+1).
\]
In a formula-generated DAG, \(\ell_e\) is a variable or constant.  Therefore each of the
two arguments of the quadratic gadget in \cref{lem:quadratic} has support at most one.
Replacing each quadratic occurrence by its own copy of the gadget yields three affine
factors of support at most two.

\subsection{Compiling ternary flow factors}

At a degree-three internal vertex, the flow factor is
\[
  1+z_a+z_b+z_c
\]
for its three incident edges.  Replace it by the four factors from
\cref{lem:ternary}.  After summing the fresh ternary bit, the resulting local value is
\[
  1+z_a+z_b+z_c+\kappa^{-1}z_az_bz_c.
\]
The last summand is the error term.

Among three incident edges and two sides of a vertex, two edges lie on the same side.
Choose \(a,b\) so that they share a head or share a tail.  The original selector contains
the corresponding exclusion \(1+z_az_b\).

\begin{lemma}[Pointwise ternary cancellation]\label{lem:local-cancel}
Fix an original selector assignment \(z\in\B^E\).  Let \(e_v(z)=\kappa^{-1}z_az_bz_c\) be
the error at a ternary vertex \(v\), and let \(W(z)\) be the product of all recovered
quadratic factors.  Then
\[
  e_v(z)W(z)=0.
\]
\end{lemma}

\begin{proof}
The product \(W(z)\) contains \(1+z_az_b\).  Hence the claim follows from
\cref{lem:killer} after multiplying by the remaining factors.
\end{proof}

The following elementary product lemma lets us remove all errors simultaneously.

\begin{lemma}[Product error elimination]\label{lem:prod-errors}
Let \(R\) be a commutative ring, let \(p_j,e_j,W\in R\), and suppose
\(e_jW=0\) for every \(j\) in a finite index set.  Then
\[
  \left(\prod_j(p_j+e_j)\right)W
  =\left(\prod_jp_j\right)W.
\]
\end{lemma}

\begin{proof}
Induct on the index set.  In the induction step,
\[
 (p+e)XW=pXW+eXW=pXW,
\]
because commutativity gives \(eXW=X(eW)=0\).
\end{proof}

\begin{theorem}[Graph compiler]\label{thm:graph-compiler}
Let \(G\) satisfy the degree and label conditions of \cref{lem:dag-properties}.  Over a
characteristic-two field containing \(\tau\notin\{0,1\}\), the path polynomial of \(G\)
has a support-two affine-hypercube representation with
\[
 q=2|E|+D+P,
 \qquad
 M=|V|+4D+3(|E|+P).
\]
\end{theorem}

\begin{proof}
Fix the original selector assignment \(z\).  Sum first over all fresh quadratic and
ternary gadget bits.  Because every occurrence has a distinct fresh auxiliary, repeated
finite distributivity factors these local sums.  By \cref{lem:quadratic}, the quadratic
sums recover every conditional edge factor and pair exclusion.  By
\cref{lem:ternary}, each ternary sum recovers the corresponding flow factor plus its cubic
error.  Apply \cref{lem:local-cancel,lem:prod-errors} while \(z\) is fixed.  The resulting
expression is precisely
\[
  \Sel_G(z)\prod_{e\in E}C_e(z_e).
\]
Now sum over \(z\) and invoke \cref{prop:abp-cube}.

For support, the source and sink flow factors use at most two selectors.  A non-ternary
internal flow factor uses at most two selectors, and a ternary vertex contributes a
constant placeholder together with four support-two gadget factors.  Every quadratic
gadget factor has support at most two.  The auxiliary count is
\[
 |E|+D+(|E|+P)=2|E|+D+P.
\]
The implemented factor list contains one flow placeholder for every vertex, four factors
per ternary occurrence, and three factors per quadratic occurrence, giving
\[
 |V|+4D+3(|E|+P).
\]
\end{proof}

\begin{remark}
The constant placeholder at a ternary vertex is retained in the formal representation.
Dropping such unit factors could improve numerical constants, but is irrelevant to the
class theorem and is not used in the certified accounting.
\end{remark}

\section{Sharp quantitative bounds}\label{sec:counts}

The generic bounds \(P\le 4|E|\) and \(D\le |V|\) already imply
\(q\le26s\) and \(M\le70s\).  A sharper pair-count identity improves both constants.

\begin{lemma}[Exact ordered-pair count]\label{lem:pairs}
For the formula DAG \(G_\Phi\),
\[
  P=4a.
\]
\end{lemma}

\begin{proof}
Induct on the formula.  A leaf creates no pair.  A multiplication gate takes disjoint
unions and adds a connector from a sink of outdegree zero to a source of indegree zero, so
it creates no new common-head or common-tail pair.  An addition gate creates two edges
with the fresh source as common tail, giving the two ordered pairs \((e_1,e_2)\) and
\((e_2,e_1)\), and two edges with the fresh sink as common head, giving two more ordered
pairs.  Thus every addition gate contributes exactly four.
\end{proof}

\begin{proof}[Proof of \cref{thm:finite-main}]
Apply \cref{thm:graph-compiler} to \(G_\Phi\).  By
\cref{lem:graph-counts,lem:pairs} and \(D\le|V|\),
\begin{align*}
 q
 &=2|E|+D+P\\
 &\le 2(l+4a+u)+(2l+2a)+4a\\
 &=4l+14a+2u\\
 &\le14(l+a+u)=14s.
\end{align*}
Similarly,
\begin{align*}
 M
 &=|V|+4D+3(|E|+P)\\
 &\le5|V|+3|E|+3P\\
 &=5(2l+2a)+3(l+4a+u)+12a\\
 &=13l+34a+3u\\
 &\le34(l+a+u)=34s.
\end{align*}
Value preservation and the support bound are supplied by
\cref{lem:dag-properties,thm:graph-compiler}.
\end{proof}

\section{The class-level lift}\label{sec:class}

We now prove \cref{thm:class-main}.  It is useful to keep the outer nondeterministic cube
and the compiler cube conceptually separate.

\subsection{Forward containment}

\begin{proposition}\label{prop:forward}
If \(\charac(\F)=2\) and \(\tau\in\F\setminus\{0,1\}\), then
\[
  \Ncl(\VPe(\F))\subseteq\VNPtwo(\F).
\]
\end{proposition}

\begin{proof}
Let \(f\in\Ncl(\VPe)\).  By \cref{def:nclosure}, there exist
\(g\in\VPe\) and polynomially bounded \(p,r\) such that
\[
 f_n(\mathbf{x})
 =\sum_{\mathbf{b}\in\B^{p(n)}}
   g_{r(n)}(\mathbf{x},\mathbf{b}).
\]
Choose formulas \(\Phi_m\) for \(g_m\) of size bounded by a polynomial \(s(m)\).  Apply
\cref{thm:finite-main} to \(\Phi_{r(n)}\), treating both \(\mathbf{x}\) and
\(\mathbf{b}\) as ordinary input variables.  We obtain a fresh cube
\(\mathbf{c}\in\B^{q_n}\), with \(q_n\le14s(r(n))\), and support-two factors
\(L_{n,j}(\mathbf{x},\mathbf{b},\mathbf{c})\), with
\(M_n\le34s(r(n))\), such that
\[
 g_{r(n)}(\mathbf{x},\mathbf{b})
 =\sum_{\mathbf{c}\in\B^{q_n}}
   \prod_{j=1}^{M_n}L_{n,j}(\mathbf{x},\mathbf{b},\mathbf{c}).
\]
Substitution and finite-sum flattening give
\[
 f_n(\mathbf{x})
 =\sum_{(\mathbf{b},\mathbf{c})\in\B^{p(n)}\times\B^{q_n}}
   \prod_{j=1}^{M_n}L_{n,j}(\mathbf{x},\mathbf{b},\mathbf{c}).
\]
The disjoint union of the two cubes has dimension \(p(n)+q_n\), which is polynomially
bounded.  Reindexing variables preserves the support-two property.  The number of factors
is polynomially bounded.  Hence \(f\in\VNPtwo\).
\end{proof}

The machine-checked class theorem uses the certified coarse bounds \(26s\) and \(70s\)
inside its polynomial-bound witnesses; the separately certified sharp theorem
\(14s,34s\) yields the same mathematical containment.

\subsection{Reverse containment}

\begin{proposition}\label{prop:reverse}
Over every field,
\[
  \VNPtwo(\F)\subseteq\Ncl(\VPe(\F)).
\]
\end{proposition}

\begin{proof}
Let
\[
 f_n(\mathbf{x})=
 \sum_{\mathbf{b}\in\B^{q(n)}}\prod_{j=1}^{M(n)}L_{n,j}(\mathbf{x},\mathbf{b})
\]
be a support-two representation with polynomially bounded \(q,M\).  For fixed \(n\), the
inner product has a formula of size \(O(M(n))\): each affine factor involving at most two
variables has constant formula size, and the factors can be multiplied by a binary tree or
a right-associated product.  The number of variables and total degree are polynomially
bounded, the latter by \(M(n)\).  These inner products therefore form a family in
\(\VPe\), and the displayed Boolean sum witnesses membership in \(\Ncl(\VPe)\).
\end{proof}

\begin{proof}[Proof of \cref{thm:class-main}]
A characteristic-two field different from \(\F_2\) contains an element
\(\tau\notin\{0,1\}\).  Apply \cref{prop:forward,prop:reverse}.
\end{proof}

\section{Width-one ABPs and the field classification}\label{sec:classification}

\subsection{The literature bridge}

A product \(L_1\cdots L_M\) is computed by a width-one ABP consisting of a path with
successive edge labels \(L_1,\ldots,L_M\).  If each \(L_j\) has support at most two, this
is a \(\mathrm{w}+\) width-one ABP in the notation of
\cite{BringmannIkenmeyerZuiddam2018}.  Consequently,
\[
 \VNPtwo\subseteq\VNPwplus\subseteq\VNPone.
\]
The artifact formalizes the first inclusion through a literature-shaped internal class in
which each factor has total degree at most one.  It deliberately does not identify the
project's unrestricted internal affine-list class definitionally with the published class.

For the opposite upper bound, \(\VPone\subseteq\VPe\): a polynomial-size product of affine
forms is a polynomial-size formula.  Monotonicity of nondeterministic closure gives
\[
 \VNPone=\Ncl(\VPone)
 \subseteq\Ncl(\VPe)=\VNPe.
\]
Valiant's characterization \(\VNPe=\VNP\) then yields
\[
 \VNPwplus\subseteq\VNPone\subseteq\VNP.
\]

\begin{theorem}[Characteristic-two width-one equality]\label{thm:char2-width1}
If \(\charac(\F)=2\) and \(\F\neq\F_2\), then
\[
 \VNPwplus(\F)=\VNPone(\F)=\VNP(\F).
\]
\end{theorem}

\begin{proof}
By \cref{thm:class-main} and Valiant's theorem,
\[
 \VNP(\F)=\Ncl(\VPe(\F))=\VNPtwo(\F).
\]
The bridge and upper bounds above give
\[
 \VNP
 =\VNPtwo
 \subseteq\VNPwplus
 \subseteq\VNPone
 \subseteq\VNP.
\]
All containments are therefore equalities.
\end{proof}

\begin{proof}[Proof of \cref{cor:classification}]
If \(|\F|>2\), either \(\charac(\F)\neq2\), in which case the support-two theorem of
Bringmann--Ikenmeyer--Zuiddam applies, or \(\charac(\F)=2\) and \(\F\neq\F_2\), in which
case apply \cref{thm:char2-width1}.  If \(|\F|=2\), then \(\F\cong\F_2\), and
Bringmann--Ikenmeyer--Zuiddam proved the strict inclusion
\(\VNPone(\F_2)\subsetneq\VNP(\F_2)\).  This rules out equality even with unrestricted
affine labels and hence also in the support-two subclass.
\end{proof}

\begin{remark}[What the theorem does not require]
No \(\F_4\) subfield is needed.  The only field-theoretic input is the existence of
\(\tau\notin\{0,1\}\), equivalently \(|\F|>2\) for a field.  This includes all infinite
fields of characteristic two and all finite fields \(\F_{2^m}\) with \(m\ge2\).
\end{remark}

\section{Formal verification}\label{sec:formal}

The proof was formalized in \lean\ 4 \cite{Lean4} on top of Mathlib
\cite{Mathlib2020}.  The support-two layer is append-only relative to the earlier
support-three development.  Its principal source modules are summarized in
\cref{tab:lean-map}.

\begin{longtable}{@{}p{0.29\textwidth}p{0.65\textwidth}@{}}
\caption{Correspondence between the mathematical proof and the formal artifact.}
\label{tab:lean-map}\\
\toprule
\textbf{Module / theorem} & \textbf{Mathematical role}\\
\midrule
\endfirsthead
\toprule
\textbf{Module / theorem} & \textbf{Mathematical role}\\
\midrule
\endhead
\code{TernaryIdentity.lean} &
Denominator-free and normalized versions of \cref{lem:quadratic,lem:ternary};
\code{kappa_ne_zero}; Boolean error killer; formal non-vanishing firewall.\\[2mm]
\code{SupportTwoGadget.lean} &
Explicit three-factor quadratic and four-factor ternary affine lists; product evaluation;
support bounds.\\[2mm]
\code{CompilerTwo.lean} &
Disjoint auxiliary types, finite-cube flattening, raw value with errors visible, and a
conditional theorem eliminating errors only under a pointwise hypothesis.\\[2mm]
\code{GraphCompilerTwo.lean} &
Degree-three pigeonhole choice, same-head/same-tail error cancellation, graph value theorem,
support bound, and exact counts \(q=2|E|+D+P\),
\(M=|V|+4D+3(|E|+P)\).\\[2mm]
\code{MainCompilerTwo.lean} &
Formula-level value and support theorems; coarse bounds \(q\le26s\), \(M\le70s\).\\[2mm]
\code{SharpConstants.lean} &
Machine-checked identity \(P=4a\); sharp bounds \(q\le14s\), \(M\le34s\).\\[2mm]
\code{ClassLevelTwo.lean} &
Definition of the internal support-two class; forward and reverse containments; equality
\(\Ncl(\VPe)=\VNPtwo\); finite-field and existential-\(\tau\) forms; width-one literature
bridge.\\[2mm]
\code{GadgetMinimality.lean} &
Achievement of one auxiliary and three factors, plus a limited zero-auxiliary degree lower
bound.  No full minimality theorem is claimed.\\[2mm]
\code{RegressionTestsTwo.lean} &
Positive examples and negative firewall tests, including failure at \(\tau=0,1\), failure
of auxiliary reuse, and the genuine support-three degree of an uncompiled ternary flow
factor.\\[2mm]
\code{AxiomAuditTwo.lean} &
\code{#print axioms} audit on the principal new theorems.\\
\bottomrule
\end{longtable}

\subsection{Reported build and trust boundary}

The supplied artifact report records:
\begin{itemize}[leftmargin=1.6em]
\item a successful \code{lake build} under \lean\ \code{v4.28.0} and Mathlib
\code{v4.28.0}, comprising 8064 jobs with no errors;
\item no declarations using \code{sorry}, \code{admit}, a new \code{axiom},
\code{unsafe}, \code{opaque}, \code{native_decide}, or \code{@[implemented_by]};
\item \code{#print axioms} output for 52 principal support-two theorems containing only
\code{propext}, \code{Classical.choice}, and \code{Quot.sound}.
\end{itemize}
The manuscript-production environment was able to inspect the archive statically but did
not contain a Lake installation, so it did not independently rerun the recorded build.
This distinction should be preserved in any public artifact statement until a separate
reproduction is completed.

No published theorem is assumed as a \lean\ axiom.  The machine-checked endpoint is the
internal equality \(\Ncl(\VPe)=\VNPtwo\) and its internal width-one bridge.  Valiant's
\(\VNPe=\VNP\), the characteristic-not-two theorem, and the \(\F_2\) separation are
imported only in the paper-level corollary.

\subsection{Artifact availability}

At the date of this draft, the public repository is
\begin{center}
\url{https://github.com/RexHannes/vnp1-char2-formalization}.
\end{center}
The support-two archive used for this manuscript has SHA-256 digest
\[
\code{6f6ccac3d858b1eb71493b69d7708fa8488a220cbd1039c2ee07354df2d4c139}.
\]
The public default branch should be checked against this digest or an identified support-two
commit before the paper is circulated as artifact-complete.

\section{Related work}\label{sec:related}

Valiant introduced the algebraic complexity framework and \(\VNP\)
\cite{Valiant1979,Valiant1980}.  Ben-Or and Cleve showed that polynomial-size formulas and
width-three ABPs are polynomially equivalent \cite{BenOrCleve1992}.  Allender and Wang
proved explicit limitations of exact width-two ABPs \cite{AllenderWang2016}.

Bringmann, Ikenmeyer and Zuiddam studied small-width ABPs, border complexity, and
nondeterministic closures in a unified framework
\cite{BringmannIkenmeyerZuiddam2018,BringmannIkenmeyerZuiddam2017}.  Their Theorem~4.2
identifies \(\VNPone\) with \(\VNP\) when the characteristic is not two; their support-two
refinement appears as Theorem~6.3 in the extended version.  Their Proposition~9.1 proves
strictness over \(\F_2\), while their discussion explicitly leaves all other
characteristic-two fields open.  They also note an \(\F_4\) representation of \(1+xy\),
which shows that the \(\F_2\) obstruction does not automatically extend to larger
characteristic-two fields.  The present \(\tau\)-gadget works uniformly without requiring
an \(\F_4\) subfield.

Dutta et al. proved universality results for \emph{border} width-two ABPs in
characteristic two and showed that a suitable power of a polynomial with a small formula
has a small border-width-two representation \cite{DuttaEtAl2024}.  That work concerns
approximate deterministic width two.  The present theorem concerns exact nondeterministic
width one; neither result subsumes the other.

\section{Limitations and open directions}\label{sec:open}

\paragraph{Constants.}
The certified list keeps a unit placeholder for every ternary vertex and uses ordered pair
exclusions.  Removing redundant unit factors or using a different selector may improve
\(14\) and \(34\).  Such improvements would not change the class theorem.

\paragraph{Local minimality.}
The artifact proves that the characteristic-two gadget achieves one auxiliary bit and three
factors, and certifies a limited degree lower bound in the zero-auxiliary setting.  It does
not formalize the full impossibility of a two-factor positive-auxiliary gadget.  No
optimality claim is part of the main theorem.

\paragraph{The exceptional field.}
Over \(\F_2\), unrestricted nondeterministic width one is strictly weaker than \(\VNP\).
The internal structure of support-bounded subclasses over \(\F_2\) remains a separate
question.

\paragraph{Alternative direct compilers.}
The graph selector gives a modular and machine-checked proof.  A direct activation-state
compiler on the formula tree might reduce constants and shorten the semantic argument, but
such a route should be treated as a new proof until separately formalized and audited.

\paragraph{Border complexity.}
The exact hypercube identities do not provide a deterministic width-two matrix product,
control an approximation parameter, or extract Frobenius roots.  Extending them to the
border-width-two setting would require an additional theorem, not a reinterpretation of
the present proof.

\section{Conclusion}

Over every characteristic-two field other than \(\F_2\), a binary arithmetic formula can
be compiled into a Boolean hypercube sum of a product of affine forms, each supported on at
most two variables, with linear overhead.  The construction works because a uniform
quadratic gadget and a ternary gadget can be integrated with a path selector whose local
exclusions annihilate the ternary error pointwise.  The resulting class equality fills the
characteristic-two range left open in the width-one nondeterministic programme and, together
with prior work, makes \(\F_2\) the unique exceptional field.

\appendix

\section{A more explicit selector proof}\label{app:selector}

For completeness, we isolate the combinatorial fact used in \cref{prop:selector}.  Let
\(H\) be the selected-edge subgraph.  Pair exclusions imply
\(\indeg_H(v),\outdeg_H(v)\in\{0,1\}\).  A nonzero source factor gives
\(\outdeg_H(s)=1\), and a nonzero sink factor gives \(\indeg_H(t)=1\).  At every other
vertex, the flow factor is nonzero if and only if
\(\indeg_H(v)=\outdeg_H(v)\).

Follow the unique selected outgoing edge from \(s\).  At any internal vertex reached, the
selected incoming edge forces a selected outgoing edge.  Because rank strictly increases,
the walk cannot repeat a vertex and must terminate.  It cannot terminate internally and
therefore terminates at \(t\).  Now suppose another selected edge exists.  Its component
contains neither \(s\) nor \(t\), since each of those already has its unique selected
incident edge in the constructed path.  Every vertex in the extra component has equal
selected indegree and outdegree.  Starting from any selected edge and repeatedly following
selected outgoing edges produces a directed cycle in the finite component, contradicting
strict rank increase.  Hence the selected subgraph is exactly one source--sink path.

\section{The exact compiler ledger}\label{app:ledger}

Let \(E,V,P,D\) be as in \cref{sec:compiler}.  The formal compiler uses the following
objects.

\begin{center}
\begin{tabular}{@{}lll@{}}
\toprule
Object & Count & Contribution\\
\midrule
Original edge selectors & \(|E|\) & auxiliary bits\\
Quadratic occurrences & \(|E|+P\) & one bit, three factors each\\
Ternary occurrences & \(D\) & one bit, four factors each\\
Flow placeholders & \(|V|\) & one factor each\\
\bottomrule
\end{tabular}
\end{center}

Therefore
\[
 q=|E|+(|E|+P)+D=2|E|+P+D,
\]
and
\[
 M=|V|+3(|E|+P)+4D.
\]
For formula graphs, substituting
\[
 |E|=l+4a+u,\quad |V|=2l+2a,\quad P=4a,\quad D\le|V|
\]
gives the sharp certified inequalities in \cref{sec:counts}.  The unproved proposed
identity \(D+B=2a\) is not used: the number of degree-three vertices depends on formula
shape beyond the triple \((l,a,u)\).

\section{Machine-checked scope firewalls}\label{app:firewalls}

The formal development includes negative or limiting statements designed to prevent the
main theorem from being overread.

\begin{enumerate}[leftmargin=1.8em]
\item \(abc(1+ab)\) is not the zero formal polynomial.
\item The ternary gadget without a pair exclusion retains a nonzero error.
\item \(\tau=0\) and \(\tau=1\) make \(\kappa=0\), so the normalized gadgets are illegal.
\item Sharing one fresh bit between distinct gadget occurrences changes the represented
polynomial in general.
\item A graph edge carrying an unrestricted three-variable affine label can violate the
support-two target; the theorem is stated for formula-generated graphs with one-variable
or constant labels.
\item A degree-three flow factor genuinely has support three before ternary compilation.
\item The class theorem is a representation theorem and includes no evaluator or runtime
claim.
\end{enumerate}

\begin{thebibliography}{BIZ18}

\bibitem[BIZ18]{BringmannIkenmeyerZuiddam2018}
K.~Bringmann, C.~Ikenmeyer, and J.~Zuiddam.
\newblock On algebraic branching programs of small width.
\newblock \emph{Journal of the ACM}, 65(5):32:1--32:29, 2018.
\newblock \href{https://doi.org/10.1145/3209663}{doi:10.1145/3209663}.

\bibitem[BIZ17]{BringmannIkenmeyerZuiddam2017}
K.~Bringmann, C.~Ikenmeyer, and J.~Zuiddam.
\newblock On algebraic branching programs of small width.
\newblock In \emph{32nd Computational Complexity Conference (CCC 2017)},
  volume 79 of \emph{LIPIcs}, pages 20:1--20:31, 2017.
\newblock \href{https://doi.org/10.4230/LIPIcs.CCC.2017.20}{doi:10.4230/LIPIcs.CCC.2017.20}.

\bibitem[Val79]{Valiant1979}
L.~G. Valiant.
\newblock Completeness classes in algebra.
\newblock In \emph{Proceedings of the Eleventh Annual ACM Symposium on Theory of Computing},
  pages 249--261, 1979.
\newblock \href{https://doi.org/10.1145/800135.804419}{doi:10.1145/800135.804419}.

\bibitem[Val80]{Valiant1980}
L.~G. Valiant.
\newblock Reducibility by algebraic projections.
\newblock Internal report, Department of Computer Science, University of Edinburgh, 1980.

\bibitem[BC92]{BenOrCleve1992}
M.~Ben-Or and R.~Cleve.
\newblock Computing algebraic formulas using a constant number of registers.
\newblock \emph{SIAM Journal on Computing}, 21(1):54--58, 1992.
\newblock \href{https://doi.org/10.1137/0221006}{doi:10.1137/0221006}.

\bibitem[AW16]{AllenderWang2016}
E.~Allender and F.~Wang.
\newblock On the power of algebraic branching programs of width two.
\newblock \emph{Computational Complexity}, 25(1):217--253, 2016.
\newblock \href{https://doi.org/10.1007/s00037-015-0114-7}{doi:10.1007/s00037-015-0114-7}.

\bibitem[DIK${}^{+}$24]{DuttaEtAl2024}
P.~Dutta, C.~Ikenmeyer, B.~Komarath, H.~Mittal, S.~G. Nanoti, and D.~Thakkar.
\newblock On the power of border width-2 {ABP}s over fields of characteristic 2.
\newblock In \emph{41st International Symposium on Theoretical Aspects of Computer Science
  (STACS 2024)}, volume 289 of \emph{LIPIcs}, pages 31:1--31:16, 2024.
\newblock \href{https://doi.org/10.4230/LIPIcs.STACS.2024.31}{doi:10.4230/LIPIcs.STACS.2024.31}.

\bibitem[Bür00]{Burgisser2000}
P.~Bürgisser.
\newblock \emph{Completeness and Reduction in Algebraic Complexity Theory}.
\newblock Algorithms and Computation in Mathematics, vol.~7. Springer, 2000.
\newblock \href{https://doi.org/10.1007/978-3-662-04179-6}{doi:10.1007/978-3-662-04179-6}.

\bibitem[MP08]{MalodPortier2008}
G.~Malod and N.~Portier.
\newblock Characterizing Valiant's algebraic complexity classes.
\newblock \emph{Journal of Complexity}, 24(1):16--38, 2008.
\newblock \href{https://doi.org/10.1016/j.jco.2006.09.006}{doi:10.1016/j.jco.2006.09.006}.

\bibitem[dMU21]{Lean4}
L.~de Moura and S.~Ullrich.
\newblock The {Lean 4} theorem prover and programming language.
\newblock In \emph{Automated Deduction---CADE 28}, volume 12699 of \emph{LNCS},
  pages 625--635. Springer, 2021.
\newblock \href{https://doi.org/10.1007/978-3-030-79876-5_37}{doi:10.1007/978-3-030-79876-5\_37}.

\bibitem[Com20]{Mathlib2020}
The mathlib Community.
\newblock The {Lean} mathematical library.
\newblock In \emph{Proceedings of the 9th ACM SIGPLAN International Conference on
  Certified Programs and Proofs}, pages 367--381, 2020.
\newblock \href{https://doi.org/10.1145/3372885.3373824}{doi:10.1145/3372885.3373824}.

\end{thebibliography}

\end{document}
